% --- Archon physics formalization source begin ---
% archon:physics
% archon:covers PhyXMiniProblems/problem_phyx_mini_0276.lean
% archon:source-report reports/phyx_mini/problem_phyx_mini_0276.source.json

\chapter{Physics problem phyx\_mini\_0276}
\label{ch:PhyXMiniProblems_problem_phyx_mini_0276}

\paragraph{Problem source.}
\#\# Physical scenario

A grandfather clock has a pendulum that consists of a thin brass disk of radius \$r = 15.00\$ cm and mass \$1.000\$ kg that is attached to a long thin rod of negligible mass. The pendulum swings freely about an axis perpendicular to the rod and through the end of the rod opposite the disk, as shown in figure.

\#\# Question

If the pendulum is to have a period of \$2.000\$ s for small oscillations at a place where \$g = 9.800\$ m/s\$\textasciicircum{}2\$, what must be the rod length \$L\$ to the nearest tenth of a millimeter?

\#\# Answer choices

- A: A: 0.8023 m

- B: B: 0.8197 m

- C: C: 0.8545 m

- D: D: 0.8315 m

\#\# Figure caption (auxiliary; use the image as primary evidence)

The image features a vertical rod with an axis of rotation marked at the top. The rod extends downward and is connected to a circular object, resembling a sphere or disk, at the bottom. The object is labeled with its radius \textbackslash{}( r \textbackslash{}). The distance from the axis of rotation to the center of the sphere is marked as \textbackslash{}( L \textbackslash{}). Two curved arrows around the rod indicate rotational motion. Text labels include "Rotation axis" next to the top of the rod, with labels \textbackslash{}( L \textbackslash{}) and \textbackslash{}( r \textbackslash{}) showing their respective measurements.

\#\# Dataset metadata

Resonance and Harmonics; Physical Model Grounding Reasoning, Multi-Formula Reasoning

\paragraph{Recorded answer/context.}
D: D: 0.8315 m

\paragraph{Figure/image path.}
/root/proposal\_for\_physic/science-mango/phyx\_mini\_run/phyx\_data/test\_image/276.png

\paragraph{Formalization target.}
create a compiling Lean file with sorry bodies at `PhyXMiniProblems/problem\_phyx\_mini\_0276.lean`. The Lean declarations must preserve the physical quantities, dimensions or dimensional roles, figure labels, governing-law hypotheses, and final relation expressed by this problem.
Use Mathlib/Physlib names found through LeanExplore where available. If a physics API is missing, introduce faithful local abstractions rather than scalar placeholder aliases.

\begin{theorem}[Physics formalization target]
\label{thm:physics:phyx_mini_0276:target}
\lean{PhyXMiniProblems.ProblemPhyXMini0276.grandfatherClockRodLength_inLinearizedModel_is_recordedAnswerD}
\uses{def:physics:phyx-mini-0276:phyxminiproblems-problemphyxmini0276-lengthinmeters, def:physics:phyx-mini-0276:phyxminiproblems-problemphyxmini0276-timeinseconds, def:physics:phyx-mini-0276:phyxminiproblems-problemphyxmini0276-accelerationinmeterspersecondsquared, def:physics:phyx-mini-0276:phyxminiproblems-problemphyxmini0276-grandfatherclockpendulumsetup, def:physics:phyx-mini-0276:phyxminiproblems-problemphyxmini0276-matchesprimarygrandfatherclockfigure, def:physics:phyx-mini-0276:phyxminiproblems-problemphyxmini0276-matchesgrandfatherclockpendulumscenario, def:physics:phyx-mini-0276:phyxminiproblems-problemphyxmini0276-hasgrandfatherclockpendulumproblemdata, def:physics:phyx-mini-0276:phyxminiproblems-problemphyxmini0276-satisfiesgrandfatherclockpendulumlaws, def:physics:phyx-mini-0276:phyxminiproblems-problemphyxmini0276-recordedanswerchoice, def:physics:phyx-mini-0276:phyxminiproblems-problemphyxmini0276-roundstonearesttenthmillimeter, def:physics:phyx-mini-0276:phyxminiproblems-problemphyxmini0276-isclosestdisplayedrodlength}
The assigned autoformalize agent should translate this physics problem into Lean declarations in the covered file, with theorem and lemma proof bodies written as `by sorry`.
\end{theorem}
\begin{proof}
This is an autoformalization task, not a proof task. Produce statements that can later be proved by the physics prover without weakening the physical model.
\end{proof}

% --- Archon declaration topology begin ---
\section{Declaration topology}
\begin{definition}[Mass Quantity]
  \label{def:physics:phyx-mini-0276:phyxminiproblems-problemphyxmini0276-massquantity}
  \lean{PhyXMiniProblems.ProblemPhyXMini0276.MassQuantity}
  A nonnegative physical mass
\end{definition}
\begin{proof}
  This is the declared model definition of Mass Quantity.
\end{proof}
\begin{definition}[Length Quantity]
  \label{def:physics:phyx-mini-0276:phyxminiproblems-problemphyxmini0276-lengthquantity}
  \lean{PhyXMiniProblems.ProblemPhyXMini0276.LengthQuantity}
  A nonnegative physical length
\end{definition}
\begin{proof}
  This is the declared model definition of Length Quantity.
\end{proof}
\begin{definition}[Time Quantity]
  \label{def:physics:phyx-mini-0276:phyxminiproblems-problemphyxmini0276-timequantity}
  \lean{PhyXMiniProblems.ProblemPhyXMini0276.TimeQuantity}
  A nonnegative physical time interval
\end{definition}
\begin{proof}
  This is the declared model definition of Time Quantity.
\end{proof}
\begin{definition}[Acceleration Magnitude Quantity]
  \label{def:physics:phyx-mini-0276:phyxminiproblems-problemphyxmini0276-accelerationmagnitudequantity}
  \lean{PhyXMiniProblems.ProblemPhyXMini0276.AccelerationMagnitudeQuantity}
  A nonnegative acceleration magnitude, with dimension L T⁻²
\end{definition}
\begin{proof}
  This is the declared model definition of Acceleration Magnitude Quantity.
\end{proof}
\begin{definition}[Moment Of Inertia Quantity]
  \label{def:physics:phyx-mini-0276:phyxminiproblems-problemphyxmini0276-momentofinertiaquantity}
  \lean{PhyXMiniProblems.ProblemPhyXMini0276.MomentOfInertiaQuantity}
  A scalar moment of inertia about an axis, with dimension M L²
\end{definition}
\begin{proof}
  This is the declared model definition of Moment Of Inertia Quantity.
\end{proof}
\begin{definition}[Restoring Torque Coefficient Quantity]
  \label{def:physics:phyx-mini-0276:phyxminiproblems-problemphyxmini0276-restoringtorquecoefficientquantity}
  \lean{PhyXMiniProblems.ProblemPhyXMini0276.RestoringTorqueCoefficientQuantity}
  The coefficient multiplying the angular displacement in the linearized gravitational restoring torque. Radians are dimensionless, so this has dimension M L² T⁻²
\end{definition}
\begin{proof}
  This is the declared model definition of Restoring Torque Coefficient Quantity.
\end{proof}
\begin{definition}[Torque Quantity]
  \label{def:physics:phyx-mini-0276:phyxminiproblems-problemphyxmini0276-torquequantity}
  \lean{PhyXMiniProblems.ProblemPhyXMini0276.TorqueQuantity}
  A signed torque component about the pivot axis, with dimension M L² T⁻². Unlike the restoring-coefficient magnitude, torque must allow both signs
\end{definition}
\begin{proof}
  This is the declared model definition of Torque Quantity.
\end{proof}
\begin{definition}[centimeter Unit Choices]
  \label{def:physics:phyx-mini-0276:phyxminiproblems-problemphyxmini0276-centimeterunitchoices}
  \lean{PhyXMiniProblems.ProblemPhyXMini0276.centimeterUnitChoices}
  Unit choices in which lengths are read in centimeters
\end{definition}
\begin{proof}
  This is the declared model definition of centimeter Unit Choices.
\end{proof}
\begin{definition}[mass In Kilograms]
  \label{def:physics:phyx-mini-0276:phyxminiproblems-problemphyxmini0276-massinkilograms}
  \lean{PhyXMiniProblems.ProblemPhyXMini0276.massInKilograms}
  \uses{def:physics:phyx-mini-0276:phyxminiproblems-problemphyxmini0276-massquantity}
  Read a physical mass in kilograms
\end{definition}
\begin{proof}
  This is the declared model definition of mass In Kilograms.
\end{proof}
\begin{definition}[length In Meters]
  \label{def:physics:phyx-mini-0276:phyxminiproblems-problemphyxmini0276-lengthinmeters}
  \lean{PhyXMiniProblems.ProblemPhyXMini0276.lengthInMeters}
  \uses{def:physics:phyx-mini-0276:phyxminiproblems-problemphyxmini0276-lengthquantity}
  Read a physical length in meters
\end{definition}
\begin{proof}
  This is the declared model definition of length In Meters.
\end{proof}
\begin{definition}[length In Centimeters]
  \label{def:physics:phyx-mini-0276:phyxminiproblems-problemphyxmini0276-lengthincentimeters}
  \lean{PhyXMiniProblems.ProblemPhyXMini0276.lengthInCentimeters}
  \uses{def:physics:phyx-mini-0276:phyxminiproblems-problemphyxmini0276-lengthquantity, def:physics:phyx-mini-0276:phyxminiproblems-problemphyxmini0276-centimeterunitchoices}
  Read a physical length in centimeters
\end{definition}
\begin{proof}
  This is the declared model definition of length In Centimeters.
\end{proof}
\begin{definition}[time In Seconds]
  \label{def:physics:phyx-mini-0276:phyxminiproblems-problemphyxmini0276-timeinseconds}
  \lean{PhyXMiniProblems.ProblemPhyXMini0276.timeInSeconds}
  \uses{def:physics:phyx-mini-0276:phyxminiproblems-problemphyxmini0276-timequantity}
  Read a physical time interval in seconds
\end{definition}
\begin{proof}
  This is the declared model definition of time In Seconds.
\end{proof}
\begin{definition}[acceleration In Meters Per Second Squared]
  \label{def:physics:phyx-mini-0276:phyxminiproblems-problemphyxmini0276-accelerationinmeterspersecondsquared}
  \lean{PhyXMiniProblems.ProblemPhyXMini0276.accelerationInMetersPerSecondSquared}
  \uses{def:physics:phyx-mini-0276:phyxminiproblems-problemphyxmini0276-accelerationmagnitudequantity}
  Read an acceleration magnitude in meters per second squared
\end{definition}
\begin{proof}
  This is the declared model definition of acceleration In Meters Per Second Squared.
\end{proof}
\begin{definition}[moment Of Inertia In Kilogram Meters Squared]
  \label{def:physics:phyx-mini-0276:phyxminiproblems-problemphyxmini0276-momentofinertiainkilogrammeterssquared}
  \lean{PhyXMiniProblems.ProblemPhyXMini0276.momentOfInertiaInKilogramMetersSquared}
  \uses{def:physics:phyx-mini-0276:phyxminiproblems-problemphyxmini0276-momentofinertiaquantity}
  Read a moment of inertia in kilogram meters squared
\end{definition}
\begin{proof}
  This is the declared model definition of moment Of Inertia In Kilogram Meters Squared.
\end{proof}
\begin{definition}[restoring Coefficient In Newton Meters]
  \label{def:physics:phyx-mini-0276:phyxminiproblems-problemphyxmini0276-restoringcoefficientinnewtonmeters}
  \lean{PhyXMiniProblems.ProblemPhyXMini0276.restoringCoefficientInNewtonMeters}
  \uses{def:physics:phyx-mini-0276:phyxminiproblems-problemphyxmini0276-restoringtorquecoefficientquantity}
  Read a restoring coefficient in newton meters per radian
\end{definition}
\begin{proof}
  This is the declared model definition of restoring Coefficient In Newton Meters.
\end{proof}
\begin{definition}[torque In Newton Meters]
  \label{def:physics:phyx-mini-0276:phyxminiproblems-problemphyxmini0276-torqueinnewtonmeters}
  \lean{PhyXMiniProblems.ProblemPhyXMini0276.torqueInNewtonMeters}
  \uses{def:physics:phyx-mini-0276:phyxminiproblems-problemphyxmini0276-torquequantity}
  Read a signed pivot-axis torque component in newton meters
\end{definition}
\begin{proof}
  This is the declared model definition of torque In Newton Meters.
\end{proof}
\begin{definition}[Disk Material]
  \label{def:physics:phyx-mini-0276:phyxminiproblems-problemphyxmini0276-diskmaterial}
  \lean{PhyXMiniProblems.ProblemPhyXMini0276.DiskMaterial}
  Material named for the pendulum disk
\end{definition}
\begin{proof}
  This is the declared model definition of Disk Material.
\end{proof}
\begin{definition}[Disk Geometry]
  \label{def:physics:phyx-mini-0276:phyxminiproblems-problemphyxmini0276-diskgeometry}
  \lean{PhyXMiniProblems.ProblemPhyXMini0276.DiskGeometry}
  Mass-distribution and thickness idealization of the disk
\end{definition}
\begin{proof}
  This is the declared model definition of Disk Geometry.
\end{proof}
\begin{definition}[Rod Idealization]
  \label{def:physics:phyx-mini-0276:phyxminiproblems-problemphyxmini0276-rodidealization}
  \lean{PhyXMiniProblems.ProblemPhyXMini0276.RodIdealization}
  Idealization of the suspension rod
\end{definition}
\begin{proof}
  This is the declared model definition of Rod Idealization.
\end{proof}
\begin{definition}[Pivot Condition]
  \label{def:physics:phyx-mini-0276:phyxminiproblems-problemphyxmini0276-pivotcondition}
  \lean{PhyXMiniProblems.ProblemPhyXMini0276.PivotCondition}
  Mechanical condition at the upper support
\end{definition}
\begin{proof}
  This is the declared model definition of Pivot Condition.
\end{proof}
\begin{definition}[Motion Plane]
  \label{def:physics:phyx-mini-0276:phyxminiproblems-problemphyxmini0276-motionplane}
  \lean{PhyXMiniProblems.ProblemPhyXMini0276.MotionPlane}
  Plane of the clock-pendulum motion
\end{definition}
\begin{proof}
  This is the declared model definition of Motion Plane.
\end{proof}
\begin{definition}[Oscillation Regime]
  \label{def:physics:phyx-mini-0276:phyxminiproblems-problemphyxmini0276-oscillationregime}
  \lean{PhyXMiniProblems.ProblemPhyXMini0276.OscillationRegime}
  Approximation used for the requested period
\end{definition}
\begin{proof}
  This is the declared model definition of Oscillation Regime.
\end{proof}
\begin{definition}[Figure Point]
  \label{def:physics:phyx-mini-0276:phyxminiproblems-problemphyxmini0276-figurepoint}
  \lean{PhyXMiniProblems.ProblemPhyXMini0276.FigurePoint}
  Distinguished points visible or geometrically determined in the image
\end{definition}
\begin{proof}
  This is the declared model definition of Figure Point.
\end{proof}
\begin{definition}[Figure Label]
  \label{def:physics:phyx-mini-0276:phyxminiproblems-problemphyxmini0276-figurelabel}
  \lean{PhyXMiniProblems.ProblemPhyXMini0276.FigureLabel}
  Text and mathematical labels visible in the image
\end{definition}
\begin{proof}
  This is the declared model definition of Figure Label.
\end{proof}
\begin{definition}[Rotation Direction]
  \label{def:physics:phyx-mini-0276:phyxminiproblems-problemphyxmini0276-rotationdirection}
  \lean{PhyXMiniProblems.ProblemPhyXMini0276.RotationDirection}
  The two directions indicated by the curved swing arrow
\end{definition}
\begin{proof}
  This is the declared model definition of Rotation Direction.
\end{proof}
\begin{definition}[Grandfather Clock Pendulum Figure]
  \label{def:physics:phyx-mini-0276:phyxminiproblems-problemphyxmini0276-grandfatherclockpendulumfigure}
  \lean{PhyXMiniProblems.ProblemPhyXMini0276.GrandfatherClockPendulumFigure}
  \uses{def:physics:phyx-mini-0276:phyxminiproblems-problemphyxmini0276-lengthquantity, def:physics:phyx-mini-0276:phyxminiproblems-problemphyxmini0276-figurepoint, def:physics:phyx-mini-0276:phyxminiproblems-problemphyxmini0276-figurelabel, def:physics:phyx-mini-0276:phyxminiproblems-problemphyxmini0276-rotationdirection}
  Physical readout of the one-panel figure. The separate rim attachment and disk center prevent the auxiliary caption's center-distance ambiguity from being built into the model
\end{definition}
\begin{proof}
  This is the declared model definition of Grandfather Clock Pendulum Figure.
\end{proof}
\begin{definition}[Grandfather Clock Pendulum Setup]
  \label{def:physics:phyx-mini-0276:phyxminiproblems-problemphyxmini0276-grandfatherclockpendulumsetup}
  \lean{PhyXMiniProblems.ProblemPhyXMini0276.GrandfatherClockPendulumSetup}
  \uses{def:physics:phyx-mini-0276:phyxminiproblems-problemphyxmini0276-massquantity, def:physics:phyx-mini-0276:phyxminiproblems-problemphyxmini0276-lengthquantity, def:physics:phyx-mini-0276:phyxminiproblems-problemphyxmini0276-timequantity, def:physics:phyx-mini-0276:phyxminiproblems-problemphyxmini0276-accelerationmagnitudequantity, def:physics:phyx-mini-0276:phyxminiproblems-problemphyxmini0276-momentofinertiaquantity, def:physics:phyx-mini-0276:phyxminiproblems-problemphyxmini0276-restoringtorquecoefficientquantity, def:physics:phyx-mini-0276:phyxminiproblems-problemphyxmini0276-torquequantity, def:physics:phyx-mini-0276:phyxminiproblems-problemphyxmini0276-diskmaterial, def:physics:phyx-mini-0276:phyxminiproblems-problemphyxmini0276-diskgeometry, def:physics:phyx-mini-0276:phyxminiproblems-problemphyxmini0276-rodidealization, def:physics:phyx-mini-0276:phyxminiproblems-problemphyxmini0276-pivotcondition, def:physics:phyx-mini-0276:phyxminiproblems-problemphyxmini0276-motionplane, def:physics:phyx-mini-0276:phyxminiproblems-problemphyxmini0276-oscillationregime, def:physics:phyx-mini-0276:phyxminiproblems-problemphyxmini0276-grandfatherclockpendulumfigure}
  Independent physical quantities for the apparatus. The pivot-to-center distance and both moments of inertia are stored independently and related to the printed lengths only by the figure and governing-law interfaces below. The Physlib oscillator represents the linearized angular coordinate
\end{definition}
\begin{proof}
  This is the declared model definition of Grandfather Clock Pendulum Setup.
\end{proof}
\begin{definition}[Matches Primary Grandfather Clock Figure]
  \label{def:physics:phyx-mini-0276:phyxminiproblems-problemphyxmini0276-matchesprimarygrandfatherclockfigure}
  \lean{PhyXMiniProblems.ProblemPhyXMini0276.MatchesPrimaryGrandfatherClockFigure}
  \uses{def:physics:phyx-mini-0276:phyxminiproblems-problemphyxmini0276-lengthinmeters, def:physics:phyx-mini-0276:phyxminiproblems-problemphyxmini0276-grandfatherclockpendulumsetup}
  Primary-image incidences and labels. In particular, L is the distance from the pivot to the disk's top rim, not to its center. Consequently the center-of-mass lever arm is L + r
\end{definition}
\begin{proof}
  This is the declared model definition of Matches Primary Grandfather Clock Figure.
\end{proof}
\begin{definition}[Matches Grandfather Clock Pendulum Scenario]
  \label{def:physics:phyx-mini-0276:phyxminiproblems-problemphyxmini0276-matchesgrandfatherclockpendulumscenario}
  \lean{PhyXMiniProblems.ProblemPhyXMini0276.MatchesGrandfatherClockPendulumScenario}
  \uses{def:physics:phyx-mini-0276:phyxminiproblems-problemphyxmini0276-grandfatherclockpendulumsetup}
  Qualitative apparatus data and idealizations stated in the problem
\end{definition}
\begin{proof}
  This is the declared model definition of Matches Grandfather Clock Pendulum Scenario.
\end{proof}
\begin{definition}[Has Grandfather Clock Pendulum Problem Data]
  \label{def:physics:phyx-mini-0276:phyxminiproblems-problemphyxmini0276-hasgrandfatherclockpendulumproblemdata}
  \lean{PhyXMiniProblems.ProblemPhyXMini0276.HasGrandfatherClockPendulumProblemData}
  \uses{def:physics:phyx-mini-0276:phyxminiproblems-problemphyxmini0276-massinkilograms, def:physics:phyx-mini-0276:phyxminiproblems-problemphyxmini0276-lengthinmeters, def:physics:phyx-mini-0276:phyxminiproblems-problemphyxmini0276-lengthincentimeters, def:physics:phyx-mini-0276:phyxminiproblems-problemphyxmini0276-timeinseconds, def:physics:phyx-mini-0276:phyxminiproblems-problemphyxmini0276-accelerationinmeterspersecondsquared, def:physics:phyx-mini-0276:phyxminiproblems-problemphyxmini0276-grandfatherclockpendulumsetup}
  The stated numerical data and positivity conditions. rodLength\_L remains an unknown positive physical length; no answer choice is assumed here
\end{definition}
\begin{proof}
  This is the declared model definition of Has Grandfather Clock Pendulum Problem Data.
\end{proof}
\begin{definition}[Satisfies Grandfather Clock Pendulum Laws]
  \label{def:physics:phyx-mini-0276:phyxminiproblems-problemphyxmini0276-satisfiesgrandfatherclockpendulumlaws}
  \lean{PhyXMiniProblems.ProblemPhyXMini0276.SatisfiesGrandfatherClockPendulumLaws}
  \uses{def:physics:phyx-mini-0276:phyxminiproblems-problemphyxmini0276-massinkilograms, def:physics:phyx-mini-0276:phyxminiproblems-problemphyxmini0276-lengthinmeters, def:physics:phyx-mini-0276:phyxminiproblems-problemphyxmini0276-timeinseconds, def:physics:phyx-mini-0276:phyxminiproblems-problemphyxmini0276-accelerationinmeterspersecondsquared, def:physics:phyx-mini-0276:phyxminiproblems-problemphyxmini0276-momentofinertiainkilogrammeterssquared, def:physics:phyx-mini-0276:phyxminiproblems-problemphyxmini0276-restoringcoefficientinnewtonmeters, def:physics:phyx-mini-0276:phyxminiproblems-problemphyxmini0276-torqueinnewtonmeters, def:physics:phyx-mini-0276:phyxminiproblems-problemphyxmini0276-grandfatherclockpendulumsetup}
  The rigid-body and small-oscillation laws used in the calculation: a thin uniform disk has center-axis inertia m r² / 2; the scalar parallel-axis theorem adds m d² at the pivot; the exact signed gravitational torque is -m g d sin θ; its derivative at the downward equilibrium is -m g d, expressed with Mathlib's local HasDerivAt contract; the corresponding positive linearized restoring coefficient is m g d;
\end{definition}
\begin{proof}
  This is the declared model definition of Satisfies Grandfather Clock Pendulum Laws.
\end{proof}
\begin{definition}[Answer Choice]
  \label{def:physics:phyx-mini-0276:phyxminiproblems-problemphyxmini0276-answerchoice}
  \lean{PhyXMiniProblems.ProblemPhyXMini0276.AnswerChoice}
  Labels of the four displayed rod-length choices
\end{definition}
\begin{proof}
  This is the declared model definition of Answer Choice.
\end{proof}
\begin{definition}[meters]
  \label{def:physics:phyx-mini-0276:phyxminiproblems-problemphyxmini0276-answerchoice-meters}
  \lean{PhyXMiniProblems.ProblemPhyXMini0276.AnswerChoice.meters}
  \uses{def:physics:phyx-mini-0276:phyxminiproblems-problemphyxmini0276-answerchoice}
  Rod length in meters printed beside each answer label
\end{definition}
\begin{proof}
  This is the declared model definition of meters.
\end{proof}
\begin{definition}[recorded Answer Choice]
  \label{def:physics:phyx-mini-0276:phyxminiproblems-problemphyxmini0276-recordedanswerchoice}
  \lean{PhyXMiniProblems.ProblemPhyXMini0276.recordedAnswerChoice}
  \uses{def:physics:phyx-mini-0276:phyxminiproblems-problemphyxmini0276-answerchoice}
  The answer label recorded in the supplied dataset metadata
\end{definition}
\begin{proof}
  This is the declared model definition of recorded Answer Choice.
\end{proof}
\begin{definition}[Rounds To Nearest Tenth Millimeter]
  \label{def:physics:phyx-mini-0276:phyxminiproblems-problemphyxmini0276-roundstonearesttenthmillimeter}
  \lean{PhyXMiniProblems.ProblemPhyXMini0276.RoundsToNearestTenthMillimeter}
  \uses{def:physics:phyx-mini-0276:phyxminiproblems-problemphyxmini0276-lengthquantity, def:physics:phyx-mini-0276:phyxminiproblems-problemphyxmini0276-lengthinmeters, def:physics:phyx-mini-0276:phyxminiproblems-problemphyxmini0276-answerchoice}
  Agreement after rounding meters to the nearest 0.0001 m, which is one tenth of a millimeter. Mathlib's round breaks exact half-way ties upward
\end{definition}
\begin{proof}
  This is the declared model definition of Rounds To Nearest Tenth Millimeter.
\end{proof}
\begin{definition}[Is Closest Displayed Rod Length]
  \label{def:physics:phyx-mini-0276:phyxminiproblems-problemphyxmini0276-isclosestdisplayedrodlength}
  \lean{PhyXMiniProblems.ProblemPhyXMini0276.IsClosestDisplayedRodLength}
  \uses{def:physics:phyx-mini-0276:phyxminiproblems-problemphyxmini0276-lengthquantity, def:physics:phyx-mini-0276:phyxminiproblems-problemphyxmini0276-lengthinmeters, def:physics:phyx-mini-0276:phyxminiproblems-problemphyxmini0276-answerchoice}
  A displayed choice is at least as close as every listed alternative
\end{definition}
\begin{proof}
  This is the declared model definition of Is Closest Displayed Rod Length.
\end{proof}
% --- Archon declaration topology end ---
% --- Archon physics formalization source end ---
